\section{Conclusion, Limitations and Future Work}\label{sec:conclusion}
EviMatch shows that a transparent pipeline, trained on 8{,}000 public pairs, can order candidates better
than similarity baselines when resumes and postings are genuinely unseen, while giving a reviewer the resume
sentence behind every requirement and the signals behind every score. It also shows what such a system
cannot yet do: accuracy is modest, the semantic evidence does not improve it, and redaction leaves
measurable proxy effects. The leakage analysis is of independent use, because comparisons on the official
split of this benchmark reward a model for recognising resumes.

\textbf{Limitations.}
(i) One benchmark, with undocumented label provenance, noisy labels (job-description identity explains
27\% of label variance and six of seven duplicated pairs disagree), partly anonymised resumes, and no stated
licence; conclusions may not transfer to other labelling schemes.
(ii) Absolute accuracy is low enough that outcomes must not be automated.
(iii) The skill lexicon is hand-written (350 entries) and English-only, and coverage outside technology and
business roles is thin.
(iv) Redaction is rule based and removes explicit identifiers; name-to-group assignment in the audit is
statistical.
(v) We could not evaluate an LLM-only baseline or audit one: the API key available to us had no request
quota. Statements about language-model bias therefore rest on the cited literature, not on our measurements.
(vi) The deletion test measures faithfulness to the model, not usefulness to people; no user study was run.
(vii) The artefact deployed in the prototype is trained on all 8{,}000 pairs; all reported numbers come from
held-out evaluations.

\textbf{Future work.} Concrete next steps are: (a) run the prepared LLM baseline and audit with a pinned
open-weight model and fixed prompts, so that the comparison is reproducible; (b) a study in which recruiters rate
whether the displayed evidence changes their decisions and time; (c) validation on a second corpus whose labels
come from a documented process; (d) link skills to an open ontology and report extraction coverage by
occupation; (e) a pool-normalised ranking mode that uses the other candidates of a posting; and (f) an
intersectional audit with larger samples and name lists validated against external data.

\section*{Ethics and AI-Assistance Statement}
The system is decision support for human reviewers. The benchmark contains resumes of unknown consent
status; no personal data beyond the public benchmark were collected, and the prototype does not store uploads.
Large-language-model tools were used to help write code and prose for this manuscript; all experiments were run,
and every number checked, from the stored result files, and the authors are responsible for the content.

\section*{Data and Code Availability}
Code, experiment scripts, result files with provenance, the evaluation protocol with its deviation log, and
this manuscript's sources are in the project repository (URL to be added at submission). The benchmark is
available from its Hugging Face page; it is not redistributed.
